\section{Related Work}
\label{sec:related}

\paragraph{Agentic and adaptive RAG.}
Self-RAG~\cite{selfrag2024} and CRAG~\cite{crag2024} introduced
reflective and corrective retrieval. Adaptive-RAG~\cite{adaptiverag2024}
routes queries by predicted complexity, and
Probing-RAG~\cite{probingrag2025} extends this idea with probes on the
model's internal states. Search-R1~\cite{searchr1_2025} trains agentic
retrieval with reinforcement learning, and RAG-Critic~\cite{ragcritic2025}
adds an iterative critic. None of these methods conditions routing on
hop distance in a typed graph or exposes a typed graph-lookup tool.

\paragraph{GraphRAG and multi-hop QA.}
Microsoft GraphRAG~\cite{graphrag2024} and its
successors~\cite{lightrag2024,hipporag2_2025,graphr1_2025} build
entity--relation graphs at indexing time. Han et al.~\cite{ragvsgraphrag2025}
and GraphRAG-Bench~\cite{graphragbench2026} report that GraphRAG does
not outperform vector RAG on every query type. Our RQ2 results suggest one reason such verdicts differ: whether
citations are scored on the retrieved set or on the answer. Of the multi-hop benchmarks MuSiQue~\cite{musique2022},
MultiHop-RAG~\cite{multihoprag2024}, and
2WikiMultihopQA~\cite{twowiki2020}, we use the first and last to test
RQ1 to RQ3 outside the requirements domain.

\paragraph{RAG for requirements traceability.}
LiSSA~\cite{lissa2025}, TVR~\cite{niu_tvr2025}, and Graph-RAG for
compliance checking~\cite{graphragcompliance2024} each evaluate a single
regulated domain, without stratifying queries by hop depth and with a
single judge.

\paragraph{Citation evaluation and judge reliability.}
ALCE~\cite{alce2023} defined citation precision, recall, and $F_1$ for
generated answers. Wallat et al.~\cite{wallat2024} show that ALCE-style
correctness does not imply faithfulness, which motivates the use of LLM
judges. RAGChecker~\cite{ragchecker2024} provides a single-judge
framework, and self-preference bias~\cite{selfpref_bias2024} motivates
judge ensembles. Feinstein and Cicchetti~\cite{feinstein1990} describe
the prevalence paradox of Cohen's $\kappa$, and Gwet's
AC1~\cite{gwet2008} is the correction we report alongside it. Our
analysis for RQ4 adds a quantity this literature does not track, the
agreement of a judge with its own earlier verdicts after the retrieval
state or the date changes.
